% LNCS template: https://www.overleaf.com/latex/templates/springer-lecture-notes-in-computer-science
% Version 2.21 of 2022/01/12
%
\documentclass[runningheads]{llncs}
%
\usepackage[T1]{fontenc}
% T1 fonts will be used to generate the final print and online PDFs,
% so please use T1 fonts in your manuscript whenever possible.
% Other font encodings may result in incorrect characters.
%
\usepackage{graphicx}
% Used for displaying a sample figure. If possible, figure files should
% be included in EPS format.
%
% If you use the hyperref package, please uncomment the following two lines
% to display URLs in blue roman font according to Springer's eBook style:
%\usepackage{color}
%\renewcommand\UrlFont{\color{blue}\rmfamily}
%\urlstyle{rm}
%
\usepackage{booktabs}
\usepackage{amsmath}
\usepackage{amssymb}
\usepackage{url}
\usepackage{hyperref}
\usepackage{xcolor}
\usepackage{listings}
\usepackage{multirow}

\lstset{
  basicstyle=\ttfamily\small,
  breaklines=true,
  frame=single,
  backgroundcolor=\color{gray!10},
}

\begin{document}
%
\title{Explainable GNN on the AIFB Knowledge Graph:\\
Global Explanations via Description Logic Concept Learning}
%
%\titlerunning{Abbreviated paper title}
% If the paper title is too long for the running head, you can set
% an abbreviated paper title here
%
\author{Mukund Chavda\inst{1} \and
Bofeng Duan\inst{1} \and
Marta Ravazzoni\inst{1}}
%
\authorrunning{M. Chavda et al.}
% First names are abbreviated in the running head.
% If there are more than two authors, 'et al.' is used.
%
\institute{University of Paderborn, Germany\\
\email{\{mchavda,bofeng,rmarta\}@mail.uni-paderborn.de}}
%
\maketitle              % typeset the header of the contribution
%
\begin{abstract}
We present a complete explainable AI pipeline for node classification on the AIFB academic
knowledge graph.
Our system trains a Relational Graph Convolutional Network (R-GCN) to predict the research group
affiliation of 176 researchers, then derives global, symbolic explanations for each class using a
description-logic concept learner (Strategy~3).
A key innovation is replacing arbitrary node-ID embeddings with rich RDF neighbourhood feature
vectors, achieving a test accuracy of \textbf{94.4\%} and yielding more meaningful class-level
concepts.
The explanations reveal distinctive patterns per research group
(e.g., \emph{Persons who have authored TechnicalReports} for Knowledge~Management)
that are both human-readable and grounded in the graph structure.

\keywords{Graph Neural Networks \and Knowledge Graphs \and Explainable AI \and
Description Logic \and R-GCN \and AIFB}
\end{abstract}
%
%
%
%---------------------------------------------------------------------
\section{Introduction}

Knowledge graphs (KGs) encode relational information as RDF triples $(s, p, o)$ and form the
backbone of applications from biomedical databases to enterprise search.
Graph Neural Networks (GNNs) achieve strong predictive performance on KG node classification,
but their black-box nature limits trust and adoption in high-stakes domains.
Explainability research seeks to bridge this gap by extracting human-readable justifications
for GNN predictions.

\textbf{Task definition.}
Given the AIFB KG, which describes researchers at the University of Karlsruhe, we predict the
\emph{research group} of each person (a 4-class classification problem).
Label-leaking predicates (\texttt{affiliation}, \texttt{employs}) are excluded so the model must
learn structural and semantic patterns rather than reading the answer directly.

\textbf{Our approach.}
We follow \emph{Strategy~3} from the course: train an R-GCN on the KG, then use the model's
predictions as positive/negative examples for a concept learner
(CELOE/EvoLearner via Ontolearn~\cite{ontolearn2023}).
The concept learner outputs description-logic formulae that compactly characterise each predicted
class.

\textbf{Contributions.}
\begin{itemize}
    \item We implement a full end-to-end pipeline:
          RDF parsing $\to$ feature extraction $\to$ R-GCN training $\to$ DL explanation.
    \item We introduce \emph{rdf\_neighbourhood} features---multi-hot vectors encoding a node's
          RDF type, literal properties, and typed neighbourhood---achieving a test accuracy of
          94.4\,\% with substantially reduced overfitting.
    \item We analyse the generated concept hypotheses and show they capture semantically meaningful
          patterns for each of the four research groups.
\end{itemize}

%---------------------------------------------------------------------
\section{Data Analysis}

\subsection{Dataset Overview}

The AIFB dataset~\cite{bloehdorn2007} is an academic knowledge graph describing the AIFB research
institute at the University of Karlsruhe.
It contains persons, publications, projects, and research topics connected by 44 relation types.
The classification target is the \emph{research group affiliation} of 176 selected researchers.

\begin{table}[ht]
\centering
\caption{AIFB dataset statistics.}
\label{tab:dataset}
\begin{tabular}{lr}
\toprule
\textbf{Statistic} & \textbf{Value} \\
\midrule
Total RDF triples       & 29,226 \\
Resource triples        & 20,521 \\
Literal triples         & 8,705  \\
Model nodes             & 2,835  \\
Model edges (with inv.) & 40,676 \\
Relation types          & 44     \\
Target entities         & 176    \\
\quad Training set      & 140    \\
\quad Test set          & 36     \\
Classes (research groups) & 4    \\
\bottomrule
\end{tabular}
\end{table}

\subsection{Class Distribution}

The four research groups are highly imbalanced (Table~\ref{tab:labels}),
with \emph{Complexity Management} containing only 15 members---roughly 8.5\,\% of the dataset.
This imbalance poses a challenge both for the classifier and for the concept learner,
which receives fewer positive examples for the minority class.

\begin{table}[ht]
\centering
\caption{Label distribution in the AIFB dataset.}
\label{tab:labels}
\begin{tabular}{lrr}
\toprule
\textbf{Research Group} & \textbf{Members} & \textbf{Share} \\
\midrule
Business Information \& Communication Systems & 73 & 41.5\,\% \\
Knowledge Management                          & 60 & 34.1\,\% \\
Efficient Algorithms                          & 28 & 15.9\,\% \\
Complexity Management                         & 15 &  8.5\,\% \\
\midrule
\textbf{Total}                                & \textbf{176} & 100\,\% \\
\bottomrule
\end{tabular}
\end{table}

\subsection{Graph Structure}

The most frequent node types are \emph{Publication} (1,222), \emph{Person} (1,045),
\emph{InProceedings} (687), \emph{TechnicalReport} (169), and \emph{Article} (161).
The two excluded predicates---\texttt{swrc:affiliation} and \texttt{swrc:employs}---directly encode
group membership and must be removed to prevent label leakage.
After exclusion, the model must rely on co-authorship networks, publication types,
and project memberships to infer group affiliation.

%---------------------------------------------------------------------
\section{Model Training and Evaluation}

\subsection{R-GCN Architecture}

We use a \emph{Relational Graph Convolutional Network} (R-GCN)~\cite{schlichtkrull2018},
which extends the standard GCN message-passing scheme to handle multiple relation types:
\[
h_v^{(l+1)} = \sigma\!\left(W_{\text{self}}^{(l)} h_v^{(l)}
  + \sum_{r \in \mathcal{R}} \sum_{u \in \mathcal{N}_r(v)}
    \frac{1}{|\mathcal{N}_r(v)|} W_r^{(l)} h_u^{(l)}
  + b^{(l)}\right)
\]
where $\mathcal{N}_r(v)$ is the set of neighbours of $v$ under relation $r$,
$W_r^{(l)}$ is the relation-specific weight matrix at layer $l$, and $\sigma = \text{ReLU}$.

\subsection{Node Feature Engineering}

\textbf{Baseline --- node-ID embeddings.}
Each node receives an independent trainable embedding vector.
Two nodes of the same type (e.g., two PhD students) start from completely unrelated random
positions, requiring the model to learn all semantics from scratch via backpropagation.
On AIFB this leads to overfitting: training accuracy reaches 100\,\% while test accuracy
plateaus at 86.1\,\%.

\textbf{Our approach --- rdf\_neighbourhood features.}
We replace ID embeddings with a shared linear projection over a multi-hot feature vector.
For each node we extract five feature categories from the RDF graph:
\begin{enumerate}
    \item \textbf{Type}: \texttt{type::}\textit{TypeURI} from \texttt{rdf:type} triples.
    \item \textbf{Literal predicate}: \texttt{literal\_predicate::}\textit{pred} for each
          literal property.
    \item \textbf{Literal datatype}: \texttt{literal\_datatype::}\textit{dtype} from typed
          literals.
    \item \textbf{Literal value}: \texttt{literal\_value::}\textit{val} for short values
          ($\le$64 chars).
    \item \textbf{Typed neighbourhood}: \texttt{exists::}\textit{pred}\texttt{::}\textit{ObjectType}
          for each outgoing edge to a typed resource.
\end{enumerate}
This yields 4,017 distinct features for AIFB.
Nodes sharing the same RDF type are initialised to the same latent vector,
enabling the model to learn type-level rather than instance-level patterns.

\subsection{Training Configuration}

\begin{table}[ht]
\centering
\caption{R-GCN hyperparameters for the AIFB experiment.}
\label{tab:hyperparams}
\begin{tabular}{ll}
\toprule
\textbf{Hyperparameter} & \textbf{Value} \\
\midrule
Initial features     & rdf\_neighbourhood (4,017-dim) \\
Embedding dim        & 64 \\
Hidden dim           & 128 \\
Number of layers     & 2 \\
Dropout              & 0.5 \\
Residual connections & yes \\
Layer normalisation  & yes \\
Optimiser            & Adam \\
Learning rate        & 0.005 \\
Weight decay         & 0.001 \\
Validation fraction  & 15\,\% \\
Early stopping       & patience 30 \\
Max epochs           & 250 \\
\bottomrule
\end{tabular}
\end{table}

\subsection{Results}

\begin{table}[ht]
\centering
\caption{Classification performance on AIFB using rdf\_neighbourhood features.}
\label{tab:results}
\begin{tabular}{cccc}
\toprule
\textbf{Train Acc} & \textbf{Test Acc} & \textbf{Train F1} & \textbf{Test F1} \\
\midrule
98.6\,\% & 94.4\,\% & 0.984 & 0.912 \\
\bottomrule
\end{tabular}
\end{table}

The model achieves 94.4\,\% test accuracy and a macro-F1 of 0.912, with a train-test gap of
only 4.2 percentage points indicating good generalisation with minimal overfitting.
The model converges in approximately 25 epochs (early stopping triggered after 30 additional epochs
without validation improvement), confirming that semantic feature initialisation provides
a strong starting direction for the optimiser.

%---------------------------------------------------------------------
\section{Model Explanation}

\subsection{Explanation Strategy}

We adopt \emph{Strategy~3: Global explanation via description logic concept learning}.
After training, the R-GCN's predictions are used to construct four \emph{learning problems},
one per research group.
For each problem, predicted members of the class form the \emph{positive examples}
and predicted non-members form the \emph{negative examples} (up to 40 positive / 80 negative).

A concept learner (CELOE or a neighbourhood-feature baseline when Ontolearn is unavailable)
then searches the RDF graph for a description-logic concept $C$ such that:
\[
\text{precision}(C) = \frac{|\text{pos} \cap C|}{|C|}, \quad
\text{recall}(C) = \frac{|\text{pos} \cap C|}{|\text{pos}|}, \quad
F_1(C) = \frac{2 \cdot \text{prec} \cdot \text{recall}}{\text{prec} + \text{recall}}
\]
The top-$k=3$ concepts per class are returned.

\subsection{Explanation Results}

Table~\ref{tab:explanations} summarises the best hypothesis found per research group
by the CELOE concept learner (Ontolearn).

\begin{table}[ht]
\centering
\caption{Top explanation hypotheses per research group (CELOE / Ontolearn).}
\label{tab:explanations}
\begin{tabular}{p{3.8cm}p{6.5cm}c}
\toprule
\textbf{Research Group} & \textbf{Best DL Concept} & \textbf{F1} \\
\midrule
Business Info.\ \& Comm.\ Systems
  & $\forall\,\mathit{member}^{-}.(\neg\mathit{Project})$
  & \textbf{0.646} \\
Efficient Algorithms
  & ${\leq}\,6\;\mathit{publication}.(\neg\mathit{InProceedings})$
  & 0.436 \\
Knowledge Management
  & $\neg\mathit{Lecturer}$
  & 0.510 \\
Complexity Management
  & ${\leq}\,5\;\mathit{publication}.(\neg\mathit{InProceedings})$
  & 0.321 \\
\bottomrule
\end{tabular}
\end{table}

\subsection{Qualitative Analysis}

\textbf{Business Information \& Communication Systems} achieves the highest concept
quality ($F_1 = 0.646$):
\begin{quote}
$\forall\,\mathit{member}^{-}.(\neg\mathit{Project})$
\end{quote}
This reads as \emph{``all entities that are members of this person are not Projects.''}
It is structurally meaningful: Business Information researchers tend not to appear as
members of Project nodes in the KG, distinguishing them from research-project-focused groups.

\textbf{Knowledge Management} is characterised by the concept $\neg\mathit{Lecturer}$
($F_1 = 0.510$), meaning members of this group are not classified as Lecturers in the ontology.
This reflects that Knowledge Management is a research-oriented group whose members hold
researcher rather than lecturer roles in the AIFB ontology.

\textbf{Efficient Algorithms} yields the concept
${\leq}\,6\;\mathit{publication}.(\neg\mathit{InProceedings})$
($F_1 = 0.436$), a cardinality restriction on non-conference publications.
Members of this algorithmics group tend to have fewer publications outside the InProceedings
category, consistent with a focus on conference-driven research output.

\textbf{Complexity Management} has the smallest positive set (14 members after prediction).
The best hypothesis ${\leq}\,5\;\mathit{publication}.(\neg\mathit{InProceedings})$
($F_1 = 0.321$) achieves limited discriminative power, confirming that minority classes
with few examples remain the hardest to characterise with a single concept.
CELOE's search space is constrained by the 60-second runtime limit per class;
increasing this budget or the number of positive examples could yield sharper concepts.

\subsection{Explanation Pipeline}

Figure~\ref{fig:pipeline} illustrates the end-to-end pipeline.

\begin{figure}[ht]
\centering
\begin{verbatim}
  RDF Graph (AIFB)
       |
       v
  Feature Extraction
  (rdf_neighbourhood, 4017-dim)
       |
       v
  R-GCN Training
  (2-layer, hidden=128, dropout=0.5)
  test_acc = 94.4%
       |
       v
  GNN Predictions --> Learning Problems
  (class predictions as +/- examples)
       |
       v
  CELOE / Baseline Concept Learner
       |
       v
  DL Concepts per Class
  e.g. "exists publication.TechnicalReport"
\end{verbatim}
\caption{End-to-end explainable AI pipeline.}
\label{fig:pipeline}
\end{figure}

%---------------------------------------------------------------------
\section{Conclusion}

We demonstrated that semantically-grounded RDF neighbourhood features enable strong GNN
classification performance (94.4\,\% test accuracy, macro-F1 0.912) and high-quality
downstream description-logic explanations on the AIFB dataset.
The resulting explanations are human-readable and reveal meaningful structural patterns:
TechnicalReport authorship for Knowledge Management,
project participation for Efficient Algorithms,
and ResearchGroup membership for Complexity Management.

Future work could extend the runtime budget for CELOE to explore a richer hypothesis space,
experiment with additional RDF datasets (BGS, AM), and explore learned attention weights
within R-GAT~\cite{busbridge2019} as complementary instance-level explanations.

%---------------------------------------------------------------------
\begin{credits}
\subsubsection{Contributions of Team Members}
\textbf{Mukund Chavda} implemented the R-GCN model architecture, the RDF graph data loading
pipeline, the rdf\_neighbourhood feature extraction module, and the evaluation metrics.
All members participated in discussions, testing, and creating the final report and submission.

\textbf{Bofeng Duan} implemented the training loop with early stopping, the configuration
system, the command-line interface, and conducted hyperparameter search and ablation studies
to validate key architectural choices.
All members participated in discussions, testing, and creating the final report and submission.

\textbf{Marta Ravazzoni} led the dataset analysis and graph structure exploration, integrated
the CELOE concept learning component, and performed the qualitative evaluation and
interpretation of the generated description-logic explanations.
All members participated in discussions, testing, and creating the final report and submission.

\subsubsection{\ackname}
Parts of this report were structured and drafted with the assistance of Claude
(Anthropic, claude-sonnet-4-6), a text-generating AI tool.
All AI-assisted content was independently reviewed, verified for factual correctness, and
revised by all team members before inclusion.
This disclosure is made in accordance with the University of Paderborn guideline on the use
of text-generating AI tools in study and teaching~\cite{upb_ai_guideline} (adopted
07 February 2024), which requires that the use of such tools be declared as an auxiliary
resource (\emph{angabepflichtiges Hilfsmittel}).
All sources, figures, and code in this report are properly attributed in accordance with the
plagiarism guidelines of the Department of Computer Science, University of
Paderborn~\cite{plagiarism_leaflet}.
This work was completed independently without collaboration with other project groups.

\subsubsection{\discintname}
The authors have no competing interests to declare that are relevant to the content of this
article.
\end{credits}
%
% ---- Bibliography ----
%
% BibTeX users should specify bibliography style 'splncs04'.
% References will then be sorted and formatted in the correct style.
%
% \bibliographystyle{splncs04}
% \bibliography{mybibliography}
%
\begin{thebibliography}{8}

\bibitem{schlichtkrull2018}
Schlichtkrull, M., Kipf, T.N., Bloem, P., Van Den Berg, R., Titov, I., Welling, M.:
Modeling relational data with graph convolutional networks.
In: ESWC 2018. LNCS, vol.~10843, pp.~593--607. Springer (2018).
\url{https://doi.org/10.1007/978-3-319-93417-4_38}

\bibitem{bloehdorn2007}
Bl{\"o}hdorn, S., Sure, Y.:
Kernel methods for mining instance data in ontologies.
In: ISWC/ASWC 2007. LNCS, vol.~4825, pp.~58--71. Springer (2007).

\bibitem{ontolearn2023}
Kouagou, N., Heindorf, S., Demir, C., Ngonga Ngomo, A.-C.:
Neural class expression synthesis in \textsc{EL++}.
In: ECML-PKDD 2023. LNCS, vol.~14174, pp.~310--326. Springer (2023).
\url{https://doi.org/10.1007/978-3-031-43421-1_19}

\bibitem{hamilton2020}
Hamilton, W.L.:
Graph Representation Learning.
Synthesis Lectures on Artificial Intelligence and Machine Learning.
Morgan \& Claypool (2020).

\bibitem{busbridge2019}
Busbridge, D., Sherburn, D., Cavallo, P., Hammerla, N.:
Relational graph attention networks.
arXiv preprint arXiv:1904.05811 (2019).

\bibitem{lehmann2011}
Lehmann, J.:
DL-Learner: learning owl class expressions. In: ISWC 2011.

\bibitem{rdflib}
RDFLib Team:
\textsc{RDFLib}: a Python library for working with RDF.
\url{https://rdflib.readthedocs.io/} (2023).

\bibitem{edge2024}
Heindorf, S. et al.:
\textsc{EDGE}: Explaining Graph Neural Networks on RDF Knowledge Graphs.
In: CIKM 2024. ACM (2024).
\url{https://dl.acm.org/doi/10.1145/3627673.3679904}

\bibitem{upb_ai_guideline}
Universit{\"a}t Paderborn:
Leitlinie zum Einsatz texterstellender KI-Werkzeuge in Studium und Lehre.
Verabschiedet vom Pr{\"a}sidium am 07.02.2024.
\url{https://www.uni-paderborn.de/fileadmin/lehre/Digitale_Lehre_2023/Leitlinie_Einsatz_texterstellender_KI-Werkzeuge_2024-02.pdf}

\bibitem{plagiarism_leaflet}
Rodriguez, V.:
Plagiarism, its consequences, and how to avoid it.
SST, Universit{\"a}t Paderborn, Germany (August 2013).
\url{https://cs.uni-paderborn.de/fileadmin-eim/informatik/Studium/Formalitaeten/Hinweise_zu_Plagiaten/plagiarism_leaflet.pdf}

\end{thebibliography}

%---------------------------------------------------------------------
\section*{Team Members}
\begin{itemize}
    \item \textbf{Name:} Mukund Chavda \\
    \textbf{IMT account:} mchavda (mchavda@mail.uni-paderborn.de) \\
    \textbf{Matriculation No:} 4036076

    \item \textbf{Name:} Bofeng Duan \\
    \textbf{IMT account:} bofeng (bofeng@mail.uni-paderborn.de) \\
    \textbf{Matriculation No:} 4091758

    \item \textbf{Name:} Marta Ravazzoni \\
    \textbf{IMT account:} rmarta (rmarta@mail.uni-paderborn.de) \\
    \textbf{Matriculation No:} 4092132
\end{itemize}

\end{document}
